\begin{textsample}{POS Dim 1 – es – Score 37.00 – es/es\_em\_15.txt}  \label{ex:f1_pos_135}
A ÁREA DE CIÊNCIAS DA NATUREZA E SUAS TECNOLOGIAS O Currículo do Espírito Santo mantém em sua estrutura a área de conhecimento Ciências da Natureza e suas Tecnologias, como disposto na BNCC e não mais associada à Área de Matemática e suas Tecnologias, como \textbf{configurada} anteriormente nos Parâmetros Curriculares Nacionais ( PCN ) de 1999, onde eram tratadas como \textbf{uma} área única.
A área de conhecimento Ciências da Natureza está representada por três componentes curriculares : Biologia, Física e Química, \textbf{que} de modo geral aparecem no Ensino Médio como algo novo, o \textbf{que} muitas vezes causa estranheza.
No entanto, todos esses componentes estiveram presentes e articulados entre si ao longo do Ensino Fundamental e na vida do estudante, porém apareciam de forma muito fragmentada e descontextualizada, como os conteúdos de Biologia, ou em \textbf{um} único momento, como ocorrente na Química e na Física do 9{\EmojiFont °} ano, observado nos PCN e expresso no currículo do Espírito Santo ( 2009 ).
A partir do novo currículo, elaborado no Espírito Santo para o Ensino Fundamental ( 2018 ), podemos observar profundas alterações na distribuição dos objetos de conhecimento, os quais refletem diretamente no Ensino Médio, haja vista \textbf{que} a Física e a Química passam a ter espaço \textbf{que} merecem durante todo o percurso formativo do Ensino Fundamental, tendo seu início ainda no quinto ano, com o primeiro contato com Química, seguido da Física nos anos seguintes.
Durante todo o percurso esses componentes são devidamente trabalhados de acordo com as habilidades apresentadas, sendo aprofundados, consolidados e ampliados posteriormente no Ensino Médio.
Compreendendo \textbf{que} na área de Ciências da Natureza do Ensino Médio, o ensino precisa se tornar mais instigante, esse novo currículo foi elaborado com propósito de se repensar a aprendizagem de forma integrada e \textbf{que} possa contribuir para a formação integral do estudante \textbf{que} está na etapa final da educação básica, \textbf{que} possui anseios, sonhos e \textbf{que} precisa ser ouvido.
Assim, nessa nova perspectiva, nasce \textbf{um} currículo tendo em si \textbf{uma} forma de desenvolver os conhecimentos a partir das habilidades numa perspectiva mais dinâmica e autônoma, possibilitando ao estudante ampliar a compreensão sobre a vida, o nosso planeta e o Universo, bem como a capacidade de refletir, \textbf{argumentar} e propor soluções e enfrentar desafios pessoais e coletivos.
Considerando o caráter interdisciplinar da área de conhecimento, há três competências específicas \textbf{que} a BNCC apresenta para a área das Ciências da Natureza, a saber : Além das mudanças na organização dos componentes curriculares, foi necessário “ reaproximar ” ou melhor, “ reconectar ” a ciência ao universo dos jovens, \textbf{uma} vez \textbf{que} a sua presença sempre foi \textbf{uma} constante, como na tecnologia de seus celulares, na forma como as informações chegam e partem, em sua alimentação e no comunicar -se, seja de forma verbal ou gestual, e na sua \textbf{percepção} do universo, assim como a sua inserção no mesmo.
Essa necessidade de reconexão é \textbf{um} grande desafio, sinalizado em propostas anteriores de organização curricular, a exemplo do PCN ( 1999 ).
No entanto, deparamos -nos com dificuldades como a digitalização das informações, onde basta digitar \textbf{uma} única palavra e as informações mais diversas surgem diante de seus olhos.
Se as informações já estão disponíveis o \textbf{que} nos resta para ensinar?
É neste ponto \textbf{que} nos encontramos!
Do \textbf{que} adianta essa gama de informações se não houver a apropriação adequada?
Como trazer a realidade por trás das informações sem embasamento científico?
Partimos então dos princípios básicos do \textbf{método} científico, onde iniciamos com a observação do fato, ou seja, das informações já disponíveis, fazendo do digital \textbf{um} aliado e não mais \textbf{um} vilão ao invés de apenas aceitá -las como verdade, passamos para a segunda etapa, a de questionar.
A utilização dos princípios básicos da metodologia científica ainda possibilita o desenvolvimento das competências socioemocionais propostas pela BNCC, a exemplo : da responsabilidade, determinação, organização, foco, persistência, iniciativa social, assertividade, entusiasmo, empatia, respeito, confiança, tolerância ao estresse, tolerância à frustação, autoconfiança, curiosidade para aprender, imaginação criativa e interesse artístico, \textbf{que} são todas essenciais para o desenvolvimento de pesquisas.
Nesse sentido, é \textbf{preciso} \textbf{despertar} no jovem o \textbf{pesquisador} \textbf{que} ele foi \textbf{enquanto} criança, só \textbf{que} agora ele não será apenas capaz de buscar as respostas e sim criar hipóteses, testá -las, discuti -las, e desse modo, chegar às suas próprias conclusões, com os argumentos e a linguagem necessária para defendê -las de forma coerente e bem fundamentada.
Ao pensar nessa nova forma de \textbf{abordagem} para a área de Ciências da Natureza, na etapa do Ensino Médio, não se apresenta \textbf{uma} lista de conteúdo a ser seguido durante o percurso geral de formação.
Nesse sentido, os dois Campos Temáticos sugeridos na BNCC, Matéria e Energia ; Vida, Terra e Cosmo ( BRASIL, 2018, p. 549 ) foram reorganizados, sendo propostos três Campos Temáticos para a área de Ciências da Natureza capixaba : Matéria e Energia ; Vida e Evolução ; Terra e Universo, mantendo a mesma proposta do Ensino Fundamental.
Vale citar \textbf{que} busca -se \textbf{uma} perspectiva histórica na educação básica, neste caso no ensino médio, perpassando por esses três campos temáticos do currículo, na tentativa de promover compreensão ampliada do mundo por meio de conexões estabelecidas entre história da ciência, conteúdos das ciências da natureza e história da \textbf{humanidade} ( CHASSOT, 2011 ; BELTRAN, 2013 ; BELTRAN, SAITO e TRINDADE, 2014 ).
MATÉRIA E ENERGIA Ao estudar o Campo Temático Matéria e Energia, é importante compreender \textbf{que} matéria é tudo ao seu redor e energia é a capacidade de provocar \textbf{uma} mudança ou realizar trabalho.
Matéria é tudo \textbf{que} tem massa e ocupa espaço e tudo no universo é feito de matéria e energia.
Incluem -se neste campo de estudo os aspectos históricos da ciência, como a lei de conservação das massas ou lei da conservação da matéria postulada por Antoine Laurent Lavoisier em 1785, mostrando a contribuição de cada pensador na evolução cognitiva de compreensão do fenômeno.
Nesse caso, Joseph Louis Proust ( 1754-1826 ) sequencialmente formulou em 1801 a “ Lei das Proporções Constantes ”.
Em 1837, Karl Friedrich Mohr \textbf{talvez} tenha feito as primeiras declarações sobre aspectos gerais do princípio da conservação de energia, na tentativa de compreender como elementos químicos podem ser transformados em energia.
Entretanto, foi na segunda metade do \textbf{século} XIX, durante a II Revolução Industrial, \textbf{que} avanços consideráveis ocorreram nos estudos sobre conservação de energia.
Nessa teia de estudos, é \textbf{preciso} conhecer alguns personagens como Sadi Carnot ( 1796-1832 ), William Thomson ( Lord Kelvin, ( 1824-1907 ), James Prescott Joule ( 1818-1889 ), Willard Gibbs(1809-1803), Rudolf Clausius ( 1822-1888 ), entre outros. > “ Na natureza nada se cria, nada se \textbf{perde}, tudo se transforma ”. > Antoine Laurent Lavoisier ( 1743-1794 ) >"A \textbf{prova} mais convincente da conversão de calor em trabalho foi derivada de minhas experiências com o motor eletromagnético, \textbf{uma} máquina composta por ímãs e barras de ferro iniciadas por \textbf{uma} bateria elétrica". > James Prescott Joule ( 1818-1889 ) VIDA E EVOLUÇÃO Ao estudar Vida, do latim “ vita ”, busca -se compreender o estado de atividade incessante comum aos seres organizados, bem como o período \textbf{que} decorre entre o nascimento e a morte.
Por outro \textbf{lado}, ao pensar na Evolução, busca -se compreender a evolução biológica, genética ou orgânica, \textbf{que} é a mudança das características hereditárias de \textbf{uma} população de seres vivos de \textbf{uma} geração para outra.
Esse processo faz com \textbf{que} as populações de organismos mudem e se diversifiquem ao longo do tempo.
O termo"evolução"pode referir -se à evidência observacional \textbf{que} constitui o fato científico intrínseco à \textbf{teoria} da evolução biológica, ou, em acepção completa, à \textbf{teoria} em sua completude.
É fato \textbf{que} esse estudo terá interseções com a \textbf{teoria} evolutiva de Charles Darwin ( 1809-1882 ), naturalista inglês, cuja \textbf{teoria} é a base da \textbf{moderna} \textbf{teoria} sintética : a \textbf{teoria} da seleção natural.
Segundo Darwin, organismos melhor adaptados ao meio têm maiores chances de sobrevivência do \textbf{que} os menos adaptados, deixando \textbf{um} número maior de descendentes.
Os organismos mais bem adaptados são, portanto, selecionados para aquele ambiente. \textbf{Uma} \textbf{teoria} científica é por definição \textbf{um} conjunto indissociável de todas as evidências verificáveis \textbf{conhecidas} e das ideias testáveis e testadas àquelas atreladas. > “ Na história da \textbf{humanidade} ( e dos animais também ), aqueles \textbf{que} aprenderam a colaborar e improvisar foram os \textbf{que} prevaleceram ”. > Charles Darwin ( 1809-1882 ) TERRA E UNIVERSO Ao estudar a Terra e Universo é \textbf{preciso} perceber \textbf{que} até cerca de 500 anos atrás, alguns acreditavam \textbf{que} a Terra era plana e chata como \textbf{um} disco.
No \textbf{século} XVI, surgiram importantes conclusões de \textbf{que} na realidade a Terra tem a forma arredondada e \textbf{que} gira em torno do sol e em torno de si mesma, mas a confirmação disso só veio na segunda metade do \textbf{século} XX.
Nos últimos anos tem sido cada vez mais frequente a necessidade de se estudar a Terra, Amazônia, Ártico, oceanos, recursos esgotáveis, produção de alimentos, fonte de água potável, tratamento de efluente, reaproveitamento de água, matrizes energéticas, tecnologias sustentáveis, bem como estudar aspectos extraterrestres ( Universo ), como o sistema solar, planetas habitáveis, novos sistemas e galáxias, entre outros aspectos.
Percebe -se a grande importância para promover o sentido de pertencimento em \textbf{uma} visão de todo, a fim de desenvolver o respeito por todo ser vivo e ambiente.
Mas, é importante considerar, de forma crítica, diversos problemas socioambientais em \textbf{que} \textbf{vivemos} como poluição, fome, doenças, desmatamento, entre outros.
Dessa forma as competências aplicadas meio de ciências da natureza, são \textbf{urgentes} para a qualidade de vida no planeta.
Pensando numa perspectiva da educação sociocultural e crítico-humanizadora, a educação é \textbf{vista} como instrumento de transformação \textbf{inspirado} na esperança e constituído na construção de saberes, nos diálogos e práticas \textbf{que} favorecem a reinvenção de pessoas capazes de pensar a si mesmas e as realidades, significando assim, as suas existências.
Para Freire ( 2011 ), “ a existência humana não pode ser muda, silenciosa nem tampouco pode nutrir -se de falsas palavras, mas de palavras verdadeiras, com \textbf{que} os \textbf{homens} transformam o mundo.
Existir humanamente é pronunciar o mundo, é modificá -lo ”.
Ao pensar no ensino de Ciências da Natureza no Estado do Espírito Santo, pode -se \textbf{agregar} as peculiaridades da Terra : a relação com o mar, a relação com diferentes culturas, o semiárido, os arranjos produtivos característicos como o café, entre outros casos.
É papel deste currículo permitir \textbf{que} o professor leve para a sala de aula questões locais e regionais \textbf{que} se conectem com os conteúdos programáticos.
O ensino de ciências deve comparecer com \textbf{teorias} e práticas já testadas e solidificadas, contendo aspectos históricos e políticos do processo de seu desenvolvimento, capaz de mostrar ao jovem \textbf{que}, por meio da ciência, é possível se alcançar aspectos inovadores e revolucionários para o bem da sociedade.
Dessa forma, busca -se promover a divulgação científica de personagens da história da ciência, articulada à \textbf{teoria}, como nos casos de Thomas Edson e Santos Dumont, com suas realizações tecnológicas ( e.g. lâmpada elétrica incandescente, dirigíveis, motor portátil ) e aprimoramento de tecnologias, como foi o caso do telefone, chuveiro quente, microfone, aparelho de raios X, sistemas de construção pré-moldado, dirigíveis com sistemas de navegação, planador, avião, aspectos da engenharia aeronáutica, entre outros casos.
Nesse sentido, busca -se estimular o interesse e a curiosidade científica induzidos pela aprendizagem cooperativa e colaborativa ( ELIAS, 1997 ; GUIMARÃES, 2018 ; FRAGELLI, 2018 ), no processo de compreensão das relações entre ciência, tecnologia, sociedade e ambiente ( \textbf{SANTOS} e MORTIMER, 2000 ; ANGOTTI e AUTH, 2001 ; \textbf{SANTOS} e AULER, 2011 ; AIKENHEAD, 2009 ; AULER, 2007 ).
Devido a situações emergenciais no ambiente, é proeminente abordar as relações de ciência, tecnologia, sociedade \textbf{agregando} questões ambientais no ensino médio, perpassando por temáticas controvérsias e transversais de educação ambiental e situações socioculturais, inclusive com aproximações da perspectiva humanística de Paulo Freire ( PINHEIRO, SILVEIRA e BAZZO, 2007 ; LUZ, QUEIROZ e PRUDÊNCIO, 2019 ; REIS e GALVÃO, 2008 ; \textbf{SANTOS}, 2008 ).
No \textbf{que} \textbf{diz} respeito às práticas pedagógicas realizadas em sala de aula, como forma de articulação do processo de \textbf{ensino-aprendizagem} com vistas a alcançar \textbf{uma} aprendizagem significativa ( MOREIRA, 2019 ; MOREIRA e MASINI, 2001 ), é importante \textbf{agregar} às práticas algumas metodologias \textbf{contemporâneas}.
São exemplos de metodologias \textbf{contemporâneas}, as metodologias ativas, sala invertida, aprendizagem \textbf{baseada} em projetos ( ou problemas ), sequências didáticas, \textbf{abordagem} lúdicas, aula de campo, jogos pedagógicos, entre outras, a fim de enriquecer e criar \textbf{um} \textbf{espírito} científico nos estudantes ( DELIZOICOV, ANGOTTI e PERNAMBUCO, 2018 ; BENDER, 2014 ; BACICH, TANZI NETO e TREVISANI, 2015 ; BORDENAVE e PEREIRA, 2015 ; POZO e CRESPO, 2009 ).
No ensino das Ciências da Natureza, vale destacar a importância do ensino de Ciências por investigação, o qual permite construir conhecimento de forma discursiva, a partir de \textbf{problematizações} de situações experimentais ou não, perpassando por etapas de hipóteses, testes de hipóteses, investigações, anotações, organização e seriação de dados, finalizando com conclusões \textbf{argumentadas} ( CARVALHO, 2013 ).
Por fim, cada vez mais as tecnologias de informação e comunicação estão presentes na sala de aula, razão pela qual não se pode refutá -las.
No \textbf{que} \textbf{diz} respeito a essa temática, há três principais aspectos a serem levados em consideração : o primeiro \textbf{diz} respeito ao uso de celular como instrumento de registro e coleta de dados – fotografias, vídeos e áudios de relatos em aulas experimentais e aulas de campo ; o segundo \textbf{diz} respeito ao uso de aplicativos e jogos digitais, com fins pedagógicos e por fim, o uso de redes sociais em sala de aula como forma de articulação dos espaços de aula com os virtuais ( BENDER, 2014 ; FRAGELLI, 2018 ; BACICH, TANZI NETO e TREVISANI, 2015 ).
Ao pensar em \textbf{uma} educação para o futuro, é \textbf{preciso} pensar em \textbf{categorias} \textbf{que} devem estar presentes na sala de aula.
São elas : cidadania, planetaridade, sustentabilidade, virtualidade, globalização, \textbf{transdisciplinaridade}, dialogicidade/dialeticidade ( GADOTTI, 2000 ).

% matched lemmas: abordagem, agregar, argumentar, basear, categoria, configurar, conhecido, contemporâneo, despertar, dizer, enquanto, ensino-aprendizagem, espírito, homem, humanidade, inspirar, lado, moderno, método, percepção, perder, pesquisador, preciso, problematização, prova, que, santo, século, talvez, teoria, transdisciplinaridade, um, urgente, ver, vivar
\end{textsample}
